%% ============================================================
%% 分片 02：原文第 2 节「${\EUN}=4$ Super Yang-Mills Theory And $S$-Duality」（原文 glshort.tex 起于第 642 行）
%% 由 tools/harvmac2tex.py 从 arXiv 源码（hep-th/0604151v3）自动生成。
%% 只译英文散文；公式、编号、\label、\cite 一律不改。
%% ============================================================

\section{${\EUN}=4$ 超对称杨-米尔斯理论与 $S$-对偶}

\label{sec:review}

\begin{keytakeaways}
\begin{itemize}[leftmargin=1.4em,itemsep=2pt,topsep=2pt]
\item \textbf{物理图像}：十维超对称杨-米尔斯理论经维数约化给出四维 ${\EUN}=4$ 理论——十维规范场分解为四维规范场与六个标量场，十维的 16 分量旋量分解为四组四维 Weyl 旋量，即四个超荷；这是四维可能的最大超对称。
\item \textbf{数学工具}：Clifford 代数与旋量表示、$\Gamma$ 矩阵的手征性与 $\bar\Gamma$ 分次、超荷代数 $\{Q_a,Q_b\}=\sum_I\Gamma^I_{ab}P_I$、$Spin(6)\cong SU(4)$ 的 $R$-对称、复耦合 $\tau=\theta/2\pi+4\pi i/e^2$。
\item \textbf{关键结论}：四维作用量在要求关于导数二次时由两个参数 $e$ 与 $\theta$ 唯一确定；$\theta$ 是角变量（$\theta\cong\theta+2\pi$），$\tau\to\tau+1$ 是初等对称性；超荷代数在边界条件下出现附加的中心荷项，为下一小节的 $S$-对偶与几何 Langlands 纲领埋下伏笔。
\item \textbf{边界与展望}：本节只做「回顾」，真正的舞台是下一小节：$S$-对偶把 $\tau\to-1/\tau$ 与规范群的对偶 ${}^L\negthinspace G$ 联系起来。
\end{itemize}
\end{keytakeaways}


 本节回顾 ${\EUN}=4$ 超对称杨-米尔斯理论及其 $S$-对偶的若干性质。

\subsection{${\EUN}=4$ 超对称杨-米尔斯理论回顾}
\label{sec:reviewsym}

如同原始文献 \cite{ref42} 中那样，${\EUN}=4$ 超对称杨-米尔斯理论
最容易由十维的维数约化构造。根据 Nahm 定理 \cite{ref43}，十维是
超对称杨-米尔斯理论可能的最大维数；并且对给定的规范群 $G$，存在一个
非常简单的超对称拉格朗日量，而且如果我们要求拉格朗日量关于曲率是
二次的，那么它在选定几个耦合参数之外是唯一的。本文中我们总假定 $G$
紧，并记其复化为 $G_{\Bbb{C}}$。这一点与几何 Langlands 纲领的大多数
其他论述不同。

  
  
  

  

\par\medskip{\it 时空约定}

 

 我们先描述一些约定。我们将在 Lorentz 号差 $-++\dots +$ 或欧几里得
号差 $+++\dots +$ 下工作。基本上，当强调拓扑场论的问题时，我们用
欧几里得号差；而当要强调这些构造在物理上合理的酉量子场论中是自然的
时候，我们用 Lorentz 号差。

 我们把十维 Minkowski 空间 $\Bbb{R}^{1,9}$ 或欧几里得空间
$\Bbb{R}^{10}$ 的度规写成
$ds^2=\sum_{I,J=0}^9g_{IJ}\,dx^I\,dx^J=\mp (dx^0)^2+(dx^1)^2+\dots
+(dx^9)^2$。它们的对称群是 $SO(1,9)$ 或 $SO(10)$，其旋量表示
${\EUS}^+$ 与 ${\EUS}^-$ 的秩为 16。在 Lorentz 号差下它们是实表示（且
互为对偶），而在欧几里得号差下它们是复共轭表示。

 $\Gamma$ 矩阵 $\Gamma_I,\,I=0,\dots,9$（在 Lorentz 号差下为实）反转
手征性，把 ${\EUS}^\pm$ 映到 ${\EUS}^\mp$，并满足 Clifford 代数
$\{\Gamma_I,\Gamma_J\}=2g_{IJ}$。此外，算子
% 式 (2.1)
\setcounter{equation}{0}
\begin{equation}\label{eq:thop}
\bar\Gamma=\Gamma_0\Gamma_1\cdots\Gamma_9
\end{equation}
 作用在
${\EUS}^+$ 或 ${\EUS}^-$ 上为乘以 1 或 $-1$。由于
${\EUS}^+$ 与 ${\EUS}^-$ 互为对偶，我们可以不把 $\Gamma$ 矩阵看作从
${\EUS}^\pm$ 到 ${\EUS}^{\mp}$ 的映射，而看作双线性映射
$\Gamma_I:{\EUS}^+\otimes {\EUS}^+\to \Bbb{R}$ 或
$\Gamma_I:{\EUS}^-\otimes{\EUS}^-\to \Bbb{R}$。（本段中我们假定
Lorentz 号差；在欧几里得号差下，所有这类映射都映到 $\Bbb{C}$，因为
旋量是复的。）若 $\beta,\gamma\in {\EUS}^+$，习惯上把双线性映射
$\Gamma_I(\beta,\gamma)$ 写成 $\bar\beta
\Gamma_I\gamma$，它在分量下也可写成
$\sum_{a,b=1}^{16}\Gamma_{I\,ab}\beta^a\gamma^b$。（$\bar \beta$ 中的
横杠是旋量的习惯记号，或许应读作转置，而非复共轭。）一个标准的约定
是：若指标 $I_1,I_2,\dots,I_k$ 不两两不同，则定义
$\Gamma_{I_1I_2\dots I_k}$ 为零；否则它等于乘积
$\Gamma_{I_1}\Gamma_{I_2}\cdots \Gamma_{I_k}$。于是当 $k$ 为奇数时
$\Gamma_{I_1I_2\dots I_k}$ 反转手征性，并且同样可以看作双线性映射
${\EUS}^+\otimes {\EUS}^+\to \Bbb{R}$（或 ${\EUS}^-\otimes {\EUS}^-\to \Bbb{R}$）。
这些映射对 $k=1,5,9$ 是对称的，对 $k=3,7$ 是反对称的。对
$\beta,\gamma\in {\EUS}^+$，我们再次把这些双线性映射
$\Gamma_{I_1\dots I_k}(\beta,\gamma)$ 写成
$\bar\beta\Gamma_{I_1\dots I_k}\gamma$。当 $k$ 为偶数时，我们有
$\Gamma_{I_1I_2\dots I_k}:{\EUS}^\pm\to{\EUS}^\pm$，等价地说，我们有
双线性映射 $\Gamma_{I_1I_2\dots
I_k}:{\EUS}^+\otimes {\EUS}^-\to \Bbb{R}$，同样记作
$\bar\beta\Gamma_{I_1\dots I_k}\gamma$。

\par\medskip{\it 场、变换律与拉格朗日量}

 十维超对称杨-米尔斯理论的场，是规范场 $A$（即 $G$-丛 $E$ 上的一个
联络）与一个费米场 $\lambda$，后者是 ${\EUS}^+\otimes {\rm
ad}(E)$ 的一个截面；换言之，$\lambda$ 是一个取伴随表示值的正手征
旋量场。在 Lorentz 号差下 $\lambda$ 是实的，因为丛 ${\EUS}^+$ 与
${\EUS}^-$ 是实丛。在欧几里得号差下 $\lambda$ 不是实的，但它的复共轭
并不进入形式体系。协变导数为 $D=d+A$，$A$ 的曲率为
$F=D^2=dA+A\wedge A$。

 我们考虑 $A$ 与 $\lambda$ 取值于 $G$ 的实 Lie 代数，其结构常数为实。
这意味着，在 $G$ 的酉表示中，$A$ 与 $\lambda$ 取值于反厄米矩阵。这与
通常的物理约定相反，但与数学文献一致。（把场取为反厄米，对 $G=U(1)$
而言可能显得不自然——这正是通常物理约定的由来——但对一般的 $G$，
它避免了不自然的 $i$ 因子。）在一套标准的物理约定中（见 \cite{ref44} 第 4 页），协变导数为 $D=d-iA'$，其中规范场
$A'$ 是厄米的。我们所用的反厄米 $A$ 与这个厄米规范场 $A'$ 之间的关系
因此是
% 式 (2.2)
\setcounter{equation}{1}
\begin{equation}\label{eq:boloko}
A=-iA'.
\end{equation}
  $A'$ 的曲率定义为
$F'=i(D)^2=dA'-i A'\wedge A'$，所以 % 式 (2.3)
\setcounter{equation}{2}
\begin{equation}\label{eq:obolko}
F=-iF'.
\end{equation}

 我们用符号「$\Tr$」表示 $G$ 的 Lie 代数上一个不变且负定的二次型。
由于我们假定 $G$ 单，这样的二次型在相差一个常数倍的意义下唯一；我们
把它归一化，使得对 $M=S^4$，特征类 ${1\over 8\pi^2}\int_M\Tr\, F\wedge F$
取任意整数值。对 $G=SU(N)$，这个二次型可由 $N$ 维表示中的迹得到，
这正是把它记作 $\Tr$ 的缘由。由于我们的 Lie 代数由反厄米矩阵生成，
$\Tr$ 是{\it 负}定的。

 在作此理解之后，十维超对称杨-米尔斯理论的超对称变换律与拉格朗日量
可以如下描述。超对称的生成元是一个取值于 ${\EUS}^+$ 的常旋量
$\epsilon$，因而满足 % 式 (2.4)
\setcounter{equation}{3}
\begin{equation}\label{eq:goggo}
\bar\Gamma\epsilon=\epsilon.
\end{equation}
（$\bar
\Gamma$ 的定义见 \eqref{eq:thop}；我们取 $\epsilon$ 为玻色性的。）
由 $\epsilon$ 生成的超对称变换为
% 式 (2.5)
\setcounter{equation}{4}
\begin{equation}\label{eq:hygo}
\begin{aligned}\delta_S A_I & = i\bar\epsilon\Gamma_I\lambda
\\
                   \delta_S\lambda & = {1\over 2}
                   \Gamma^{IJ}F_{IJ}\epsilon.\\\end{aligned}
\end{equation}

 对任意场 $\Phi$，记号 $\delta_S \Phi$ 是\footnote{在
$\Bbb{Z}_2$-分次代数中，记号 $[A,B\}$ 表示 $AB-(-1)^{AB}BA$。}
$\sum_{a=1}^{16}[\epsilon^aQ_a,\Phi\}$ 的简写，其中 $Q_a$ 取值于
${\EUS}^-$，是十六个超对称。记号 $\delta_S$ 表示超对称变分。不变
作用量为\footnote{这是按 Lorentz 号差写出的，采用的是动能取正的通常
约定。对欧几里得号差，必须改变拉格朗日量的整体符号，才能使作用量的
玻色部分正定。} % 式 (2.6)
\setcounter{equation}{5}
\begin{equation}\label{eq:ygo}
{{\cmmib I}}_{10}={1\over e^2}\int
d^{10}x\,\Tr\left({1\over
2}F_{IJ}F^{IJ}-i\bar\lambda\Gamma^ID_I\lambda\right),
\end{equation}
 其中
$e$ 是任意常数，即规范耦合。超对称性的验证见 \cite{ref42}。最后，生成
这些超对称 \eqref{eq:hygo} 的守恒超流为
% 式 (2.7)
\setcounter{equation}{6}
\begin{equation}\label{eq:bygo}
J^I=\half\Tr\, \Gamma^{JK}F_{JK}\Gamma^I\lambda.
\end{equation}

 这个理论的玻色对称性不只是 $SO(1,9)$，而是「Poincaré 群」$\EUP$，
它是 $SO(1,9)$ 被 $\Bbb{R}^{1,9}$ 的「平移」群所作的扩张。这个平移群
同构于 $\Bbb{R}^{1,9}$ 本身（视作阿贝尔群），而 ${\EUP}$ 是一个扩张 % 式 (2.8)
\setcounter{equation}{7}
\begin{equation}\label{eq:ugy}
0\to
\Bbb{R}^{1,9}\to \EUP\to SO(1,9)\to 1.
\end{equation}
  $\Bbb{R}^{1,9}$ 的生成元称为「动量算子」$P_I$。守恒超荷 $Q_a$，
$a=1,\dots,16$ 所满足的代数为 % 式 (2.9)
\setcounter{equation}{8}
\begin{equation}\label{eq:realify}
\{Q_a,Q_b\}=\sum_{I=1}^{10}\Gamma^I_{ab}P_I.
\end{equation}
 此外，
$Q_a$ 与 $P_I$ 对易，并在 $SO(1,9)$ 下按 ${\EUS}^-$ 变换。

\par\medskip{\it 约化到四维}

 要约化到四维，我们只需让所有场与坐标 $x^4,\dots,x^9$ 无关。分量
$A_I$，$I=0,\dots,3$ 描述四维规范场
$A=\sum_{\mu=0}^3A_\mu dx^\mu$，而分量 $A_I$，$I\ge
4$ 变成四维标量场 $\phi_i=A_{i+4}$，
$i=0,\dots,5$。十维曲率 $F_{IJ}$ 有三类贡献：视指标
$I,J$ 落在范围 $4,\dots, 9$ 内的个数为零、一或二，我们分别得到
四维曲率 $F_{\mu\nu}$、标量场的导数
$D_\mu\phi_i$，或标量场的对易子 $[\phi_i,\phi_j]$。

 在四维中，作用量的玻色部分包含全部三类贡献，成为 % 式 (2.10)
\setcounter{equation}{9}
\begin{equation}\label{eq:loopy}
{{\cmmib
I}}_4={1\over e^2}\int d^4x\,\Tr\,\left({1\over
2}\sum_{\mu,\nu=0}^3F_{\mu\nu}F^{\mu\nu}+\sum_{\mu=0}^3\sum_{i=1}^6D_\mu\phi_i
D^\mu\phi_i +{1\over 2}\sum_{i,j=1}^6[\phi_i,\phi_j]^2\right).
\end{equation}

 再加上作用量中涉及费米子的部分（它同样可以用四维的语言写出，尽管
我们不这样做），这本质上就是具有最大可能超对称的唯一四维规范理论。
若 $G$ 单，且我们要求拉格朗日量关于导数是二次的（这是通向合理量子
理论的情形），则作用量在以下选择之外是唯一的：参数 $e$，以及在四维中
另一个度量 $G$-丛 $E$ 拓扑的参数： % 式 (2.11)
\setcounter{equation}{10}
\begin{equation}\label{eq:oopy}
{{\cmmib I}}_{\theta}=-{\theta\over
8\pi^2}\int \Tr \,F\wedge F.
\end{equation}
 最后这一项是 $\theta$ 乘以该丛的第二 Chern 类，或称瞬子数。参数
$e$ 与 $\theta$ 组合成一个复参数
% 式 (2.12)
\setcounter{equation}{11}
\begin{equation}\label{eq:ingo}
\tau={\theta\over 2\pi}+{4\pi i\over e^2}.
\end{equation}
 只要
我们在 $\Bbb{R}^4$ 上，或在任何满足 $H^2(M,\Bbb{Z})=0$ 的四维流形
$M$ 上（至于 $\pi_1(G)$ 开始起作用的推广，见 \cite{ref45,ref26}），
就有一个初等的对称性 $\tau\to\tau+1$；它表达的事实是：
$(1/8\pi^2)\int\Tr F\wedge F$ 取整数值，而在量子力学中人们只关心
作用量模去 $2\pi$ 的整数倍。等价地，$\theta$ 是一个角变量，
$\theta\cong\theta+2\pi$。

 十维中的 $SO(1,9)$（或 $SO(10)$）对称性，在维数约化到四维后变为
$SO(1,3)\times
SO(6)$（或 $SO(4)\times SO(6)$）。考虑到理论中存在旋量，真正的对称性
是 $Spin(1,9)$ 约化为 $Spin(1,3)\times Spin(6)$（或 $Spin(10)$ 约化为
$Spin(4)\times Spin(6)$）。群 $Spin(6)$ 同构于 $SU(4)$，被称为该理论的
「$R$ 对称群」。我们将称它为 $SU(4)_{{\cal R}}$。

 十维中的手征条件 $\bar\Gamma\epsilon=\epsilon$ 在四维中变为 % 式 (2.13)
\setcounter{equation}{12}
\begin{equation}\label{eq:yug}
\hat\Gamma
\Gamma'\epsilon=\epsilon,
\end{equation}
 其中
$\hat\Gamma=\Gamma_0\Gamma_1\Gamma_2\Gamma_3$ 度量 $Spin(1,3)$ 手征性，
$\Gamma'=\Gamma_4\Gamma_5\dots\Gamma_9$ 度量 $Spin(6)$ 手征性。
$\hat \Gamma$ 与 $\Gamma'$ 的特征值为 $\pm i$；\eqref{eq:yug} 意味着
$\Gamma'$ 的特征值是 $\hat \Gamma$ 的特征值的相反数。$\hat\Gamma$ 的
两个特征值区分 $Spin(1,3)$ 的两个旋量表示，而 $\Gamma'$ 的特征值
同样标记 $Spin(6)$ 的旋量表示。$Spin(1,3)$ 的复化为
$SL(2)\times SL(2)$，两个旋量表示对应于 $SL(2)\times SL(2)$ 的表示
$({\bf 2},{\bf 1})$ 与 $({\bf 1},{\bf 2})$（这里 $({\bf
2},{\bf 1})$ 是第一个 $SL(2)$ 的二维表示与第二个 $SL(2)$ 的一维平凡
表示作张量积；$({\bf 1},{\bf 2})$ 反之）。$Spin(6)$ 的旋量表示是
$SU(4)_{{\cal R}}$ 的定义四维表示及其对偶；我们记之为 ${\bf 4}$ 与
$\bar{\bf 4}$。我们选取定向，使得 $\hat \Gamma$ 在
$({\bf 2},{\bf 1})$ 与 $({\bf 1},{\bf 2})$ 上分别作用为 $i^{-1}$ 与
$i$，而 $\Gamma'$ 在 $\bar {\bf 4}$ 与 ${\bf
4}$ 上作用为 $i$ 与 $i^{-1}$。
 于是
\eqref{eq:yug} 意味着，四维超对称在 $Spin(1,3)\times Spin(6)\sim SL(2)\times SL(2)\times
Spin(6)$ 下按 % 式 (2.14)
\setcounter{equation}{13}
\begin{equation}\label{eq:jugo}
({\bf 2},{\bf 1},\bar{\bf 4})\oplus ({\bf 1},{\bf 2},{\bf 4}).
\end{equation}
 变换。

 费米场 $\lambda$ 在 $Spin(1,3)\times Spin(6)$ 下按同样的方式变换。

 如果我们把 $({\bf 2},{\bf 1},\bar{\bf 4})$ 型的超对称记作
$\bar Q_{AX}$（$A=1,2$，$X=1,\dots, 4$），类似地把
$({\bf 1},{\bf 2},{\bf  4})$ 型的超对称记作 $ Q_{\dot A }^Y$
（$\dot A=1,2$，$Y=1,\dots,4$），那么由这些超对称生成的代数就是
\eqref{eq:realify} 的约化，
% 式 (2.15)
\setcounter{equation}{14}
\begin{equation}\label{eq:ugo}
\begin{aligned}\{\bar Q_{AX}, Q_{\dot A}^
Y\}&=\delta_X{}^Y\sum_{\mu=0}^3\Gamma^\mu_{A\dot A}P_\mu\\
\{Q,Q\}=\{\bar Q,\bar Q\}&=0,\\\end{aligned}
\end{equation}
 其中四维动量算子 $P_\mu$ 现在生成 $\Bbb{R}^{1,3}$ 的平移。
在适当的边界条件下，会出现附加项 \cite{ref13}，位于 \eqref{eq:ugo} 的右端
（它们与在约化到四维时被丢弃的另外六个动量以及它们的磁对偶有关）。
这些项将在本节末尾短暂出场。

\begin{yizhu}[易错点]
按：本节的号差与符号约定必须成套地看：$ds^2=\mp(dx^0)^2+\cdots$ 中上符号对应 Lorentz 号差、下符号对应欧几里得号差；反厄米规范场约定 $A=-iA'$ 连带曲率 $F=-iF'$；欧几里得号差下还必须整体翻转拉格朗日量的符号才能让作用量的玻色部分正定。后面 BPS 方程与 Hitchin 方程的符号都建立在这些约定之上，抄公式时不要混用物理文献的厄米约定。
\end{yizhu}

\begin{yizhu}[含糊点拨]
按：「十维是超对称杨-米尔斯理论可能的最大维数」是 Nahm 定理（1978）的内容：无引力的超对称理论最高只能到十维，十一维只允许超引力。超荷必须是该维数旋量表示的元素，而十维的 Majorana--Weyl 旋量恰有 16 个分量；再往上一维，超荷代数就无法与规范场的表示内容相容。
\end{yizhu}

\begin{yizhu}[物理图像]
按：十维到四维的约化之所以恰好给出 ${\EUN}=4$，可以这样数：十维规范场 $A_I$ 的十个分量分解为四维规范场 $A_\mu$（$\mu=0,\dots,3$）与六个标量场 $\phi_i=A_{i+4}$，而十维的 16 分量旋量分解为四组四维 Weyl 旋量，正是 ${\EUN}=4$ 的四个超荷（每个超荷本身又是 $\bar Q_{AX}$ 这样的多重态）。超对称的个数由旋量分量数决定，而标量个数由「十减四」决定，两者被超对称代数锁定。
\end{yizhu}

\begin{yizhu}[易错点]
按：$Spin(6)\cong SU(4)$ 是 ${\EUN}=4$ 理论的 $R$-对称群，四个超荷属于 $\mathbf{4}$ 与 $\bar{\mathbf{4}}$（见式 \eqref{eq:jugo}）。$R$-对称是超荷之间的对称性，与规范群 $G$ 无关；后文讨论拓扑扭化时，正是从 $SU(4)_{{\cal R}}$ 里挑出一个 $U(1)$ 或 $SU(2)$ 子群来定义扭化，因此这里的表示分解不是装饰，而是后文构造的出发点。
\end{yizhu}

\begin{yizhu}[易错点]
按：手征性在十维与四维的记号不同：十维的 $\bar\Gamma=\Gamma_0\cdots\Gamma_9$ 特征值为 $\pm1$，而四维的 $\hat\Gamma$、$\Gamma'$ 特征值成了 $\pm i$（欧几里得号差下的约定）。式 \eqref{eq:yug} 说的不是「两者各自为正」，而是「两者的特征值互为相反数」。若按 Lorentz 号差的实特征值去套，会差出 $i$ 因子。
\end{yizhu}

\begin{yizhu}[物理图像]
按：复耦合 $\tau=\theta/2\pi+4\pi i/e^2$ 中的 $4\pi$ 不是笔误：在本文的归一化下瞬子的作用量为 $8\pi^2|k|/e^2$，于是第 $k$ 个瞬子扇区的权重可写成 $e^{2\pi i\tau k}=e^{i\theta k-8\pi^2k/e^2}$，这正是 $\theta$ 角与瞬子数配对的方式。也正因为 $\theta$ 是角变量，$\tau\to\tau+1$ 才是理论的一个初等对称性；下文的 $SL(2,\Bbb{Z})$ 对偶群便由 $T:\tau\to\tau+1$ 与 $S:\tau\to-1/\tau$ 生成。
\end{yizhu}

\begin{yizhu}[含糊点拨]
按：本段末尾预告的「附加项」就是 ${\EUN}=4$ 超荷代数在非平凡边界条件下的中心荷：被约化丢弃的六个动量 $P_{i+4}$ 的期望值（标量场的渐近值）会作为中心荷出现，而它们在 $S$-对偶下的像就是磁荷。后文正是用这些中心荷来定义 't Hooft 线算子，并进而把 Wilson/'t Hooft 算子解释为 Hecke 算子。
\end{yizhu}
